% =====================================================================
%  Abbildung 5.x  --  Gain
% =====================================================================
\begin{figure}
\centering
\begin{tikzpicture}
  \node[fui] (sl)  at (1.7, 0.75) {\texttt{hslider("gain")}};
  \node[fsig](in)  at (0.0,-0.75) {$x[n]$};
  \node[fb2] (mul) at (4.6, 0.0)  {$*$};
  \node[fsig](out) at (6.5, 0.0)  {$y[n]$};

  \draw[fw] (sl.east)  -- (3.80, 0.75) |- \pinA{mul};
  \draw[fw] (in.east)  -- (3.45,-0.75) |- \pinB{mul};
  \draw[fw] (mul.east) -- (out);

  \begin{scope}[on background layer]
    \node[fgrp, fit=(sl)(in)(mul)] (proc) {};
  \end{scope}
  \fgname{proc}{process}
\end{tikzpicture}
\caption[Faust-Blockdiagramm der Verst\"arkungsstufe]%
{Verst\"arkungsstufe (\texttt{gain.dsp}). Beide Eing\"ange der
Multiplikation liegen auf der linken Blockkante: das Audiosignal und das
Bedienelement. Der Verst\"arkungsfaktor ist damit ein Signal und keine
\"Ubersetzungszeitkonstante; er wird je Aufruf von \texttt{compute()} aus dem
Objekt gelesen. Zwischen aufeinanderfolgenden Samples besteht keine
Abh\"angigkeit.}
\label{fig:faust-gain}
\end{figure}
